\documentclass[10pt,a4paper]{amsart}
\usepackage[margin=19mm]{geometry}
\usepackage{amsmath,amssymb,mathtools}
\usepackage[scr=rsfs,cal=euler]{mathalfa}
\usepackage{xfrac}
\usepackage[unicode,hidelinks,bookmarks=false]{hyperref}
\newcommand{\ie}{i.e.\ }
\newcommand{\cf}{cf.\ }
\newcommand{\resp}{respectively\ }
\newcommand{\Subsection}{\subsection}
\providecommand{\nobreakdash}{\nobreak\hbox{-}}
\setlength{\emergencystretch}{2em}
\multlinegap0pt
\usepackage{bm}
\usepackage{bbm}
\usepackage{dsfont}
\usepackage{xspace}
\usepackage{cancel}
\usepackage{mathtools}
\usepackage{amsthm}
\makeatletter
\@ifundefined{Theorem}{\newtheorem{Theorem}{Theorem}}{}
\@ifundefined{Corollary}{\newtheorem{Corollary}[Theorem]{Corollary}}{}
\@ifundefined{Lemma}{\newtheorem{Lemma}[Theorem]{Lemma}}{}
\@ifundefined{Proposition}{\newtheorem{Proposition}[Theorem]{Proposition}}{}
\@ifundefined{Remark}{\newtheorem{Remark}[Theorem]{Remark}}{}
\makeatother
\newcommand{\D}{{\widetilde{D}_{n+2}}}
\newcommand{\I}{\mathrm{I}}
\newcommand{\TTC}{\mathrm{CC}}
\newcommand{\car}{\mathrm{K}}
\newcommand{\fc}{{\mathrm{FC}}}
\newcommand{\ld}{{\mathrm{D}_{\mathrm{L}}}}
\newcommand{\rd}{{\mathrm{D}_{\mathrm{R}}}}
\newcommand{\red}{\rightsquigarrow}
\newcommand{\sn}{{s_\circ}}
\newcommand{\s}{{s_{\bullet}}}
\newcommand{\ti}{{t_{\bullet}}}
\newcommand{\tn}{{t_\circ}}
\title[Star and weak star irreducible fully commutative elements in]{Star and weak star irreducible fully commutative elements in Coxeter groups of affine types $\widetilde{B}$ and $\widetilde{D}$}
\author{Riccardo Biagioli, Luca Costantini, Elisa Sasso}
\date{Source: arXiv:2508.08388v1 (CC BY 4.0)}
\begin{document}
\maketitle

\noindent\emph{Source and licence.} Star and weak star irreducible fully commutative elements in Coxeter groups of affine types $\widetilde{B}$ and $\widetilde{D}$. Version arXiv:2508.08388v1 (CC BY 4.0), distributed under the Creative Commons Attribution 4.0 International licence, as recorded in the arXiv licence metadata for this exact version and shown on the abstract page https://arxiv.org/abs/2508.08388v1. Full rights notice: licenses/arxiv-2508.08388-CC-BY-4.0.txt.\par\smallskip

\noindent\textit{Editorial scope.} Counted unit: the Proposition whose source label is \texttt{nuovodue} in the article (source0.tex lines 1293--1301, source label \texttt{nuovodue}), delivered together with the article's own complete original proof of it (source0.tex lines 1303--1322, one \texttt{proof} environment of 20 lines). No mathematics is written, repaired or re-derived: statement and proof are the article's own text line for line. Required context retained uncounted: eq:mst (lines 168--169); oss (lines 344--358); theorem:starope (lines 406--414); corollary:corclassification (lines 781--784); remark:tlalgebra (lines 1274--1278); nuovoprel (lines 1282--1285). Every definition, display and label of those lines is retained unchanged. Omitted from this package and not redistributed: 1--167, 170--343, 359--405, 415--780, 785--1273, 1279--1281, 1286--1292, 1302--1302, 1323--1532. Nothing inside a retained excerpt is omitted or altered: every retained body is delivered exactly as the article's lines are written, so no omission receipt of any kind accompanies this package. Citations: the retained excerpts carry no citation command at all, so no bibliography entry of the article is transcribed and no reference is redistributed. Reference adaptation: none was needed and none was made; every reference inside the delivered unit resolves inside it, so no unresolved reference or citation is produced. Printed slips kept literal: no printed word, symbol, hypothesis or number of the counted statement or of its proof was altered. Layout and class: the article's own class is \texttt{amsart} with packages including geometry, fontenc, newlfont, inputenc, babel, graphicx, amsmath, amsthm, wrapfig, latexsym, amssymb, amsfonts, tikz, subfigure; the wrapper is the lane's pinned \texttt{amsart} editorial wrapper at 10pt on A4, which restates the theorem-like environments the retained text needs and adds the article's own macro definitions that the retained text needs (\texttt{\textbackslash D}, \texttt{\textbackslash I}, \texttt{\textbackslash TTC}, \texttt{\textbackslash car}, \texttt{\textbackslash fc}, \texttt{\textbackslash ld}, \texttt{\textbackslash rd}, \texttt{\textbackslash red}, \texttt{\textbackslash s}, \texttt{\textbackslash sn}, \texttt{\textbackslash ti}, \texttt{\textbackslash tn}), with extra wrapper packages (bm, bbm, dsfont, xspace, cancel, mathtools). The delivered numbering is the unit's own. Assets: no retained span contains a figure environment, a graphics-inclusion command, a PDF-page inclusion, a table or a TikZ picture; no image file is copied into this package, the package declares no asset dependency and no image file is redistributed with it. Licence: version arXiv:2508.08388v1 of this article is distributed under the Creative Commons Attribution 4.0 International licence, as recorded in the arXiv licence metadata for this exact version and shown on the abstract page https://arxiv.org/abs/2508.08388v1. A full rights notice is delivered at licenses/arxiv-2508.08388-CC-BY-4.0.txt. Characters: every retained excerpt is pure ASCII, so the delivered engine, XeLaTeX with its default Latin Modern text font, reports no missing character.
Let $M$ be a square symmetric matrix indexed by a finite set $S$, satisfying  $m_{s,t}=m_{t,s}\in \mathbb{N}_{>0}\cup \{\infty\}$ and $m_{s,t}=1$ if and only if $s=t$, for all $s,t\in S$. The \textit{Coxeter group} $W$ associated with the \textit{Coxeter matrix} $M$ is defined by generators $S$ and relations $(st)^{m_{s,t}}=e$ if $m_{s,t}< \infty$. These relations can be rewritten as $s^2=e$ for all $s\in S$, and \begin{equation}\label{eq:mst}    
[st]_{m_{s,t}}:=\underbrace{sts\cdots}_{m_{s,t}}=\underbrace{tst\cdots}_{m_{s,t}}=[ts]_{m_{s,t}}\end{equation} for all $s,t\in S$ such that $2\le m_{s,t}< \infty$, the latter being called \textit{braid relations}. When $m_{s,t}=2$, they are simply \textit{commutation relations} $st=ts$. 

\begin{Remark}\
\label{oss}

\begin{enumerate}
    \item[(1)] If $w\red_t sw$ (respectively, $w\red_t ws$), then $\ell(tw)>\ell(w)$ (respectively, $\ell(wt)>\ell(w)$).%, since $s\in \ld(w)$ (respectively, $s\in \rd(w)$), $m_{s,t}\ge 3$ and $w\in \fc(W)$.
    \item[(2)] We have that $w$ is left (weak) irreducible if and only if $w^{-1}$ is right (weak) irreducible. Thus, $w$ is (weak) irreducible if and only if $w^{-1}$ is also (weak) irreducible.
    \item[(3)] If $3\le m_{s,t}<\infty$, define $\xi_{s,t} := [st]_{m_{s,t}-1}$ (see \eqref{eq:mst}). We have that $w\red_t sw$ (respectively, $w\red_t ws$) is a weak reduction if and only if $\xi_{s,t}$ is a prefix (respectively, $\xi_{s,t}^{-1}$ is a suffix) of $w$.    
   % and consider the pair $I:=\{s,t\}$. We define $\xi_{s,t} \in W_I$ such that $\ell(\xi_{s,t})=m_{s,t}-1$ and $\ld(\xi_{s,t})=\{s\}$. We have that $w$ is left (respectively, right) weak star reducible by $s$ w.r.t. $t$ if and only if $\xi_{s,t}$ is a prefix (respectively, $\xi_{s,t}^{-1}$ is a suffix) of $w$. ??
    \item[(4)] Assume $m_{s,t}=3$ and $w=u_0\cdots u_m$ be its CFNF, $w\red_t sw$ if and only if $s\in supp(u_0)$, $t\in supp(u_1)$ and for every $s'\in S$ such that $m_{s',t}= 3$, $s'\notin supp(u_0)$. 
    \item[(5)] Assume $s\in \ld(w)$ (resp. $s\in \rd(w)$). If $t\in S$ such that $m_{s,t}\ge 3$ and $tw\notin \fc(W)$ (respectively, $wt\notin \fc(W)$), then $w$ admits $st$ as a prefix (respectively, $ts$ as a suffix), so $w\red_t sw$ (respectively, $w\red_t ws$) is a weak reduction. 
    \item[(6)] Assume $s\ne s'$, if $w\red_t sw$ and $w\red_{t'} s'w$ (respectively, $w\red_t ws$ and $w\red_{t'} ws'$) are (weak) reductions, then $m_{x,x'}=2$ for all $x\in \{s,t\}$ and $x'\in \{s',t'\}$, therefore $sw\red_{t'} s'sw$ (respectively, $ws \red_{t'} wss'$) is a (weak) reduction. 
    \item[(7)] % w=u's' con t'\in \rd(u') e st prefisso di u', quindi t\in \ld(su') e t'\in \rd(su') allora su's' ha come prefisso t's'.  
    Let $w\red_t sw$ and $w\red_{t'} ws'$ be (weak) reductions. Then $ws'\red_t sws'$  is a (weak) reduction if and only if $sw \red_{t'} sws'$ is a (weak) reduction. 
\end{enumerate}
\end{Remark}

\begin{Theorem}\label{theorem:starope}\
    \begin{itemize}    
        \item[(a)] If $v\in \fc(W)$ is left (respectively, right) (weak) irreducible and obtained from $w$ by a series of left (respectively, right) (weak) reductions, then $v$ is unique.
        \item[(b)] Let $w,v,v'\in \fc(W)$ with $v,v'$ (weak) irreducible. If $w$ is (weak) reducible to $v$ and $v'$ with respectively sequences \begin{align} w_0&=w\red w_1\red w_2\red \cdots \red w_l=v, \label{eq:seq1}\\
    w_0'&=w\red w_1'\red w_2'\red \cdots \red w_k'=v' \label{eq:seq2};
    \end{align} then $k=l$. 
        \item[(c)] Assume $|S|>2$. If $w$ is (weak) reducible to $v$ (weak) irreducible and $supp(v)$ is complete, then $v$ is unique. 
    \end{itemize}
\end{Theorem}

\begin{Corollary}
\label{corollary:corclassification}
    If $w\in \fc(\D)$ is reducible to $v,v'\in \I(\D)$, then $v,v'$ belong to the same family of irreducible elements. In particular, if $v,v'\notin \TTC(\D)$, then $v=v'$. 
\end{Corollary}

\begin{Remark}
\label{remark:tlalgebra}
It is not difficult to see that by applying the relations (d1)-(d3), $b_{{i_1}}\cdots b_{{i_r}}=\delta^k b_x$ with $k\ge 0$ and $x\in \fc(\D)$. In particular, $s_{i_1}\in \ld(x)$ and $s_{i_r}\in \rd(x)$.
% forse meglio metterlo come un semplice remark dato che deriva strettamente della presentazione dell'algebra tl. 
\end{Remark}

\begin{Lemma}
\label{nuovoprel}
Let $w\in \fc(\D)$ and suppose there exist $s\in S$ and $s'\in \ld(w)$ with $m_{s,s'}=3$. Then $b_sb_w=b_x$ with $x\in \fc(\D)$, $f_\bullet(w)=f_\bullet(x)$ and $f_\circ(w)=f_\circ(x)$.
\end{Lemma}

\begin{Proposition}
\label{nuovodue}
Let $w\in \fc(\D)$ be reducible to $v\in \car(\D)\cup \TTC(\D)$, and let $s\notin \ld(w)$ be such that $b_sb_w=\delta ^kb_{x}$ with $x\in \fc(\D)$. 
\begin{itemize}
    \item[(a)] If $k>0$, then $v\in \car(\D)$ and $f_\bullet(x)=f_\circ(x)=1$.
    \item[(b)] If $v\in \TTC(\D)$, then $k=0$.
\end{itemize}
%If $s\notin \ld(w)$, $b_sb_w=\delta ^kb_{x}$ with $k>0$ and , then  $f_\bullet(x)=f_\circ(x)=1$ and $v\in \car(\D)$.
\end{Proposition}

\begin{proof}

Note that (b) follows by (a). 
To prove (a) we consider a reduction sequence $w_0=w\red w_1\red w_2\red \cdots \red w_l=v$ with $v=u_0\cdots u_m$ its CFNF.
We prove the result by induction on $l$.

If $l=0$, then $w=v\in \car(\D)\cup \TTC(\D)$. If $w\in \TTC(\D)$, then $sw\in \fc(\D)$ and $b_sb_{w}=b_{sw}$. So, we have $w\in \car(\D)$ and we assume $\{\s , \sn\}\subset \ld(w)$. If $s=s_j$ with $2\le j\le n$ even, then there exists $s'\in \ld(w)$ such that $m_{s,s'}=3$, so $b_sb_w=b_x$ by Lemma $\ref{nuovoprel}$. Thus, $s\in \{\ti, \tn\}$, then $b_sb_w=\delta ^kb_x$ with $k=\lfloor m/2\rfloor -1$ and $x=s_0s_1s_3s_5\cdots s_{n-1}s_{n+1}s_{n+2}\in \TTC(\D)$ by direct computation and $f_\bullet(x)=f_\circ(x)=1$. 

Assume $l>0$. First, we assume that $w\red_p rw=:w_1$. We have that $m_{s,r}=2$, otherwise if $m_{s,r}=3$, then $b_sb_w=b_y$ with $y\in \fc(\D)$ by Lemma \ref{nuovoprel}. Therefore, 
\[\delta ^kb_{x}=b_sb_w=b_sb_{r}b_{w_1}=b_{r}b_{s}b_{w_1}\] with $k>0$ and $s\notin \ld(w_1)$.
Now we distinguish two cases. Suppose $m_{s,p}=2$. By Remark \ref{remark:tlalgebra} we have that $b_sb_{w_1}=\delta ^hb_{y}$ with $h\ge 0$, $y\in \fc(\D)$ and $p\in \ld(y)$, so \[b_rb_{s}b_{w_1}=\delta ^hb_{r}b_{y}.\]
Thus by Lemma \ref{nuovoprel}, $b_rb_y=b_{x'}$ with $x'\in \fc(\D)$ and $f_\bullet(x')=f_\bullet(y)$ and $f_\circ(x')=f_\circ(y)$. Therefore, $\delta ^kb_{x}=\delta ^hb_{x'}$, hence $h=k>0$ and $x=x'$. By inductive hypothesis on $w_1$, we have that $f_\bullet(y)=f_\circ(y)=1$ and $v\in \car(\D)$, in particular \[f_\bullet(x)=f_\bullet(y)=1 \mbox{ and }f_\circ(x)=f_\circ(y)=1.\]

On the other hand, if $m_{s,p}=3$, we have that $sw_1\notin \fc(\D)$, otherwise if $sw_1\in \fc(\D)$, then $rsw_1=sw\in \fc(\D)$ and $b_sb_w=b_{sw}$. 
Thus, by Remark \ref{oss}(5), it means that $w_1\red_s pw_1:=w'$. Therefore, $w$ is star reducible to $w'$, in particular by Theorem \ref{theorem:starope}(2) there exists a sequence of length $l$ such that\[ w\red w_1'=w_1\red w_2'=w'\red \cdots \red w'_l=v'\in \car(\D)\cup \TTC(\D)\] 
by Corollary \ref{corollary:corclassification}. Thus, \[\delta ^kb_x=b_{r}b_{s}b_{w_1}=b_{r}b_{s}b_{p}b_{w'}=b_{r}b_{w'}\] with $k>0$, since $b_{s}b_{p}b_{w'}=b_{w'}$ and $s\in \ld(w')$. Moreover, since $w=rpw'$ reduced then $r\notin \ld(w')$. Therefore, by inductive hypothesis $f_\bullet(x)=f_\circ(x)=1$ and $v'\in \car(\D)$, so we have that $v=v'$ by Corollary \ref{corollary:corclassification}.

 Now, we suppose that $w\red_p wr:=w_1$. Thus, \[\delta ^kb_x=b_sb_w=b_sb_{w_1}b_{r}\] with $k>0$. By Remark \ref{remark:tlalgebra}, $b_sb_{w_1}= \delta ^hb_y$ and since $p\in \rd(w_1)$, we have that $p\in \rd(y)$. Thus, $b_yb_r=b_{x'}$ with $x'\in \fc(\D)$ and $f_\bullet(x')=f_\bullet(y)$ and $f_\circ(x')=f_\circ(y)$ by Lemma \ref{nuovoprel}. Therefore, $\delta ^k b_x=\delta^h b_{x'}$, so $k=h>0$ and $x=x'$. By inductive hypothesis on $w_1$, we have that $f_\bullet(y)=f_\circ(y)=1$ and $v\in \car(\D)$, in particular \[f_\bullet(x)=f_\bullet(y)=1 \mbox{ and }f_\circ(x)=f_\circ(y)=1.\]

\end{proof}

\end{document}
